
\newpage
% ============================================================
%  SECTION: Stateful vs. Stateless Scaling in Kubernetes
% ============================================================

\section{Scaling Stateful and Stateless Applications in Kubernetes}
\label{sec:stateful-vs-stateless}

A fundamental distinction in distributed systems design is that between
\textit{stateless} and \textit{stateful} applications, and this distinction has
direct consequences for how workloads are deployed, managed, and scaled inside
a Kubernetes cluster. Understanding this divide is essential for motivating the
core research question of this thesis.

% ------------------------------------------------------------
\subsection{Stateless Applications and Horizontal Scalability}
\label{subsec:stateless-scaling}

A stateless application does not retain any client-specific data between
requests. Each request is self-contained and can be served by any replica of
the application without prior knowledge of previous interactions
(\cite{redhat:statefulvsstateless}). This architectural property makes stateless
services trivially scalable: because every replica is functionally identical,
the Kubernetes scheduler can start or terminate any number of them at any time
without coordination. Kubernetes \texttt{Deployments} and
\texttt{ReplicaSets} are the canonical workload primitives for stateless
applications, providing identical, interchangeable pod replicas with no notion
of stable identity or ordered lifecycle (\cite{kubernetes:statefulsets}).

Horizontal scaling of a stateless service via the Horizontal Pod Autoscaler
(HPA) is therefore straightforward: when aggregate CPU or memory utilisation
exceeds a configured threshold, the HPA controller issues a scale-out
instruction to the relevant \texttt{Deployment}, and the new replicas are
scheduled and become available without any data migration, topology
reconfiguration, or disruption to existing replicas (\cite{kubernetes:autoscaling}).
Conversely, when load decreases, replicas can be terminated immediately since
no in-replica state is lost. This symmetry between scale-out and scale-in is
what makes the reactive HPA model well-suited to stateless workloads.

% ------------------------------------------------------------
\subsection{The Challenges of Scaling Stateful Applications}
\label{subsec:stateful-challenges}

Stateful applications, by contrast, maintain persistent data or internal state
that is specific to individual instances and must survive pod restarts,
rescheduling events, and scaling operations. This category includes relational
databases, key-value stores, message queues, and consensus systems. Kubernetes
provides the \texttt{StatefulSet} workload primitive to accommodate these
requirements. Unlike a \texttt{Deployment}, a \texttt{StatefulSet} guarantees
that each pod receives a stable, unique network identity (e.g.,
\texttt{redis-0}, \texttt{redis-1}), a dedicated
\texttt{PersistentVolumeClaim} that is bound across restarts, and an ordered,
sequential lifecycle for pod creation and termination
(\cite{kubernetes:statefulsets}, \cite{portworx:statefulsets}).

These guarantees, while necessary for correctness, fundamentally complicate
autoscaling. The following challenges are particularly relevant in the context
of this thesis:

\paragraph{Ordered deployment and termination.}
\texttt{StatefulSet} pods are created and deleted in strict ordinal order.
When scaling out from $n$ to $n+1$ replicas, Kubernetes will not proceed until
the $n$-th pod is fully \textit{Ready}. This serialisation ensures application-level
invariants (such as cluster quorum) are maintained but introduces latency that
can be problematic under sudden load spikes (\cite{kubernetes:statefulsets}).
A reactive autoscaler such as the HPA, which fires only after a resource
threshold has already been breached, is therefore poorly positioned to absorb
burst traffic for stateful workloads: by the time new replicas are ordered,
scheduled, initialised, and joined to the cluster, the performance degradation
may already be severe.

\paragraph{Persistent volume provisioning.}
Each new pod added to a \texttt{StatefulSet} requires a new
\texttt{PersistentVolumeClaim} to be dynamically provisioned by the storage
backend. Depending on the underlying storage class and infrastructure, this
provisioning step can itself introduce non-trivial latency, further extending
the time-to-ready for new stateful replicas compared to stateless counterparts
(\cite{portworx:statefulsets}).

\paragraph{Cluster topology and data rebalancing.}
For distributed databases and in-memory stores that shard data across their
nodes — such as Redis Cluster — adding or removing a node is not merely a
matter of starting or stopping a process. Redis Cluster partitions its keyspace
across exactly \num{16384} hash slots, each of which is assigned to a primary
node (\cite{redis:clusterspec}). When a new node is added, the cluster operator
must \textit{reshard} a subset of hash slots from existing nodes to the new
one; when a node is removed, its slots must first be migrated away. This
resharding process involves live data movement between nodes and, if performed
reactively under resource pressure, carries a non-trivial risk of increased
latency, partial unavailability, or data inconsistency
(\cite{googlecloud:zerodowntime}).

\paragraph{Quorum and availability constraints.}
Many stateful distributed systems depend on a quorum of nodes being available
to elect a leader, acknowledge writes, or service reads. Redis Cluster requires
at least one primary per shard to be available, and Kubernetes
\texttt{PodDisruptionBudget} (PDB) resources are the mechanism used to
express minimum availability constraints during voluntary disruptions such as
autoscaling events (\cite{kubernetes:disruptions}). Aggressive scale-in
triggered by a reactive autoscaler can violate quorum constraints if PDBs are
not carefully configured, risking full cluster unavailability.

\paragraph{The HPA and stateful workloads.}
While the HPA technically supports \texttt{StatefulSets} as scaling targets,
its reactive, threshold-based model is widely considered inappropriate for
database and cache workloads (\cite{hpa:bestpractices}). The HPA operates on
\textit{lagging indicators}: it reacts to resource pressure that has already
materialised. For a stateful application such as Redis, where high memory
pressure can degrade cache hit rates and high CPU pressure can stall
replication, by the time the HPA observes a breach and a new node has
completed resharding, a significant window of degraded performance will have
already elapsed. Furthermore, the HPA's default synchronisation period of
15~seconds (\cite{keda:scalingdeployments}) means that multiple polling cycles
may pass before a scaling decision is acted upon, compounding the lag.

These challenges collectively motivate the investigation of a more
\textit{proactive}, event-driven autoscaling strategy for stateful workloads,
as explored in this thesis through the use of KEDA.
